\section{Introduction}\label{sec:introduction}

A walk in a fixed environment is a Markov chain. Averaging over an
independent environment at every lattice site creates a different
process: revisiting a site also revisits information about its transition
probabilities. The directional zero--one problem asks whether the
annealed probability of escape in one prescribed direction can lie
strictly between zero and one. We resolve this problem positively in
every dimension $d\ge3$ for iid nearest-neighbor environments with
strictly positive transition probabilities. No common positive lower
bound on those probabilities is assumed.

\subsection{The model and the theorem}

Let $d\ge3$ be an integer, let $e_1,\ldots,e_d$ be the coordinate basis,
and set
\[
 U=\{\pm e_1,\ldots,\pm e_d\},\qquad
 \mathcal P_U=\left\{p\in[0,1]^U:\sum_{e\in U}p(e)=1\right\},
 \qquad \Omega=\mathcal P_U^{\ZZ^d}.
\]
Equip $\Omega$ with its product Borel $\sigma$-field.
An environment $\omega\in\Omega$ assigns a transition vector
$\omega(x,\cdot)$ to each site $x$. We call this vector the \emph{row}
at $x$. Let $\mu$ be a probability measure on $\mathcal P_U$, and let
$\PP=\mu^{\otimes\ZZ^d}$. Thus the rows are independent and identically
distributed; entries within one row need not be independent. Assume only
\begin{equation}\label{eq:ellipticity}
 \mu\bigl(\{p\in\mathcal P_U:p(e)>0\text{ for every }e\in U\}\bigr)=1.
\end{equation}
We call this \emph{strict ellipticity}. By countability, it gives
$\omega(x,e)>0$ simultaneously for all sites and steps, almost surely.
Unlike uniform ellipticity, it permits the smallest row entry to be
arbitrarily close to zero.

The path space is $\mathcal X=(\ZZ^d)^{\NN_0}$ with coordinate
maps $X_n$ and the $\sigma$-field $\sigma(X_n:n\ge0)$. For $x\in\ZZ^d$ and a fixed environment, the
\emph{quenched law} $P_{x,\omega}$ is the Markov-chain law specified by
\[
 P_{x,\omega}(X_0=x)=1,\qquad
 P_{x,\omega}(X_{n+1}=X_n+e\mid X_0,\ldots,X_n)=\omega(X_n,e).
\]
The \emph{annealed law} and its expectation are
\[
 P_x(B)=\int P_{x,\omega}(B)\,\PP(d\omega),\qquad E_x.
\]
When both environment and path are needed, we use the joint law
$\mathsf P_x(d\omega,dX)=\PP(d\omega)P_{x,\omega}(dX)$.
The path filtration is $\mathcal F_n=\sigma(X_0,\ldots,X_n)$;
on the joint space the larger filtration
$\mathcal G_n=\sigma(\omega,X_0,\ldots,X_n)$ also reveals the entire
environment. Quenched Markov arguments use $\mathcal G_n$.
The endpoint-posterior martingales later in the proof use
$\mathcal F_n$, with the environment integrated out.
An unconditioned annealed path law will also be called the \emph{raw law}.

For a nonzero vector $\ell\in\RR^d$, define
\[
 A_\ell=\{X_n\cdot\ell\longrightarrow+\infty\}.
\]
\begin{theorem}\label{thm:main}
Let $d\ge3$, let $\PP$ be an independent and identically distributed
law of nearest-neighbor transition rows on $\ZZ^d$ satisfying
\eqref{eq:ellipticity}, and let $\ell\in\RR^d\setminus\{0\}$.
Then
\[
 P_0(A_\ell)\in\{0,1\}.
\]
\end{theorem}

The direction is any fixed real vector; no rationality assumption is
made, and no assertion of one exceptional null set uniform over all
directions is needed. No inverse-transition or logarithmic moment,
drift, balance, reversibility, or speed condition supplements
\eqref{eq:ellipticity}. In the proof, a finite first moment for a
\emph{spatial radius} is derived under the hypothesis that both signs
have positive probability. No first moment for a regeneration duration,
positive speed or ballisticity assertion is used or obtained from that
estimate.

\subsection{Prior work}

The directional question goes back to Kalikow~\cite{Kalikow1981};
see also the historical discussion in~\cite{Zerner2007}.
The classical sign-union zero--one law concerns
$P_0(A_\ell\cup A_{-\ell})$, not $P_0(A_\ell)$.
It originates in the uniformly elliptic theory of
Kalikow~\cite{Kalikow1981}; its strict-ellipticity extension is
\cite[Proposition~3]{ZernerMerkl2001}, as recalled in
\cite[Equation~(1)]{Zerner2007}. Zerner and Merkl proved the
one-sign directional law in dimension two for iid strictly elliptic
environments~\cite[Theorem~1]{ZernerMerkl2001}, and Zerner gave a shorter
proof~\cite{Zerner2007}. Slonim extended the planar iid law to bounded jumps
when almost surely there is a unique infinite communicating class reachable
from every site~\cite[Theorem~2.1]{Slonim2024}. Theorem~\ref{thm:main} establishes
the directional zero--one conjecture for nearest-neighbor walks in every
higher dimension under strict ellipticity.

The difficulty in higher dimensions is that the planar arguments force
oppositely directed paths to meet by geometry; their quantitative
counterpart must be supplied differently here. Berger proved that in
dimensions $d\ge5$ an iid uniformly elliptic walk has at most one nonzero
limiting velocity~\cite[Theorem~1.1]{Berger2008}. This allows a second
velocity equal to zero and leaves open the possibility of opposite
directional escape at zero speed. Independence is also substantive:
Bramson, Zeitouni, and Zerner constructed stationary, polynomially mixing,
uniformly elliptic environments in $d\ge3$ with positive probabilities of
linear escape in opposite directions~\cite[Theorem~3]{BramsonZeitouniZerner2006}.
Their Open Problem~3 formulates the higher-dimensional question for iid
uniformly elliptic environments.

A non-uniformly elliptic special case is provided by Dirichlet-distributed
rows. Sabot and Tournier obtained positive-probability transience in coordinate
directions with positive corresponding mean drift, using reversal identities
specific to that distribution~\cite[Theorem~1]{SabotTournier2011}.
Bouchet proved the fixed-direction zero--one law in this class for $d\ge3$
by studying an accelerated walk and an invariant law for the environment
viewed from the particle~\cite[Corollary~4]{Bouchet2013}. Tournier identified
the asymptotic direction as that of the initial drift when this drift is
nonzero~\cite[Corollary~1]{Tournier2015}. The argument below applies to arbitrary iid
strictly positive rows and uses only translations of its path pieces.

The regeneration structure of Sznitman and Zerner~\cite{SznitmanZerner1999}
is a central tool. We give the factorization and width arguments needed
below, including the distinction between real projected heights and
integer record indices. Simenhaus proved that distinct deterministic
asymptotic directions occurring with positive probability must be
opposite~\cite[Proposition~1]{Simenhaus2007}. His existence theorem for
an asymptotic direction assumes transience on a nonempty open set of
directions. Drewitz and Ram\'{i}rez obtained a corresponding two-ray
description under sign-union transience on an open set, equivalently in
$d$ linearly independent directions~\cite[Theorem~1.8]{DrewitzRamirez2010}.
Those hypotheses do not follow here from the initial assumption in one
direction. In our contradiction argument, the spatial radius estimate
supplies the two rays first; a self-contained argument then proves
their antipodality. The remaining estimates concern contacts between
two arrangements of regeneration pieces.

\subsection{Proof strategy}

Two fresh-row arguments make the proof work without a common positive
lower bound on the transition probabilities. First, finite-width slab escape
(Lemma~\ref{lem:finite-supremum}) is
proved using recurrent-site trials and first visits to new columns,
in the spirit of \cite[Lemma~4]{ZernerMerkl2001}. Second, scripts launched
from a tilted record use pairwise distinct, previously unused rows.
Their annealed cost is a product of positive row averages, rather than
a product of uniform quenched lower bounds.

Suppose, toward a contradiction, that $0<P_0(A_\ell)<1$.
A slab-exit argument and fresh-half-space regeneration first show that
the path escapes in one of the two signs almost surely.
At a strict record from which the path
never falls below its new height, the future factors from the past.
The path between successive such records is a finite \emph{regeneration
piece}. Under conditioning never to backtrack, these pieces are
independent and identically distributed. Their mean height width is
finite, even for an irrational direction.

The first main step, Proposition~\ref{prop:radius}, proves that
coexistence of the two signs also forces finite mean spatial radius of
a piece. A very large lateral
excursion would create a record in a tilted direction. From that record,
a fresh continuation escaping in the opposite sign has a probability
bounded below independently of the excursion scale. This would give a
fixed positive probability to arbitrarily large finite height maxima,
contradicting the tail of their distribution.
The radius bound gives deterministic limiting rays on the two escape
events; the rays must be opposite. We may therefore reduce to
coexistence in one coordinate direction.

We then sample two independent sequences of coordinate regeneration
pieces. Place positive pieces downward by successively anchoring their
\emph{terminal} points, and place negative pieces downward in their
ordinary chronological order. Individual path words retain their
original orientation. Heights at which both sequences have a piece
boundary divide the arrangement into independent common blocks.
A \emph{contact} is a site from which both arranged paths depart.
For each dyadic interval of common-block indices, consider whether a
contact occurs. Let $m(J)$ be the sum of these probabilities over the
first $J$ such intervals, as defined precisely in
\eqref{eq:contact-count}.

The two final steps give
\[
 m(J)=O\!\left(\frac{J}{\log J}\right),\qquad
 m(J)\ge c\,\frac{J}{\log\log J}\quad\text{for large }J,
\]
with $c>0$, which are incompatible.
For the upper bound, Proposition~\ref{prop:entropy-upper} compares the
arrangement with independent bridges of prescribed heights. Randomizing
the number of blocks keeps the cost of replacing each chunk by bridges
of its realized height bounded.
An independent stretch of blocks before the comparison interval turns the
lateral alignment cost into an increment of displacement entropy. These
increments telescope over dyadic scales, while finite mean widths make
the bridge contact probabilities summable on those scales. A
relative-entropy bound converts these two estimates into the stated bound
for the expected number of contact intervals.

For the lower bound, view a positive bridge in chronological order and
start the negative path just below its random endpoint. A union bound
over all possible endpoints would be too expensive. Instead, the
conditional probabilities of the possible endpoints form a vector of
annealed martingales. Lemma~\ref{lem:posterior-maxima} bounds the sum of
their running maxima by an exponential tail on a logarithmic scale.
In Proposition~\ref{prop:contact-bound}, first-contact prefixes are transferred to
one shared environment before they are classified by quenched exit
probabilities; the posterior functions remain in the unrevealed
annealed path filtration. This separation controls the endpoint-dependent
alignment without imposing a martingale property after the environment
is revealed. A range of dyadic height intervals without a contact forces
a late first contact near the top. Applying the estimate at the many
regeneration boundaries in that range yields many contact height scales.
The finite mean common-block width then converts the height
count into the block-index count $m(J)$.

Section~\ref{sec:regeneration} proves the path-weight and regeneration
facts. Section~\ref{sec:radius} establishes the spatial moment bound.
Section~\ref{sec:arrangement} reduces to coordinate heights and defines
the common arrangement. Sections~\ref{sec:entropy} and
\ref{sec:posterior} prove the entropy and endpoint-posterior estimates.
Section~\ref{sec:contradiction} gives the lower contact count and
completes the theorem.
